\documentclass[10pt,a4paper]{amsart}
\usepackage[margin=19mm]{geometry}
\usepackage{amsmath,amssymb,mathtools}
\usepackage[scr=rsfs,cal=euler]{mathalfa}
\usepackage{xfrac}
\usepackage[unicode,hidelinks,bookmarks=false]{hyperref}
\newcommand{\ie}{i.e.\ }
\newcommand{\cf}{cf.\ }
\newcommand{\resp}{respectively\ }
\newcommand{\Subsection}{\subsection}
\providecommand{\nobreakdash}{\nobreak\hbox{-}}
\setlength{\emergencystretch}{2em}
\multlinegap0pt
\usepackage{dsfont}
\usepackage{csquotes}
\usepackage{amssymb}
\usepackage{tikz}
\newtheorem{proposition}{Proposition}
\newcommand{\R}{\mathbb{R}}
\newcommand{\abs}[1]{|#1|}
\newcommand*\diff{\mathop{}\!\mathrm{d}}
\newcommand{\mat}[1]{
    \begin{pmatrix} #1 \end{pmatrix}
}
\newcommand{\norm}[1]{\|#1\|}
\title{Korn's inequality from the viewpoint of calculus of variations}
\author{Gabriele Cassese}
\date{Source: arXiv:2603.22431v3 (CC BY 4.0)}
\begin{document}
\maketitle

\noindent\emph{Source and licence.} Korn's inequality from the viewpoint of calculus of variations. Version arXiv:2603.22431v3 (CC BY 4.0), distributed by arXiv under the Creative Commons Attribution 4.0 International licence, http://creativecommons.org/licenses/by/4.0/, as recorded in the arXiv licence metadata for this exact version and shown on the abstract page https://arxiv.org/abs/2603.22431v3. Full licence notice: \texttt{licenses/arxiv-2603.22431-CC-BY-4.0.txt}.\par\smallskip

\noindent\textit{Editorial scope.} Counted unit: the article's proposition prop:lowerbound2 (source0.tex lines 1903--1905). The article's complete original proof of it is delivered with it (source0.tex lines 1906--2058, one \texttt{proof} environment of 153 lines). No mathematics is written, repaired or re-derived: the statement and the proof are the article's own lines, in the article's own notation, and no retained line is substituted or re-flowed.  No further source-local context is needed: every cross-reference of every retained line resolves inside the two delivered spans.  Nothing was omitted from any retained span, so the package carries no omission record.  No retained line contains a citation, so the wrapper carries no bibliography and no bibliography entry.  No retained line contains a figure environment, an image-inclusion command or drawing code, so no image is delivered and the package declares no asset dependency.  Editorial formatting only, in addition: an AMS \texttt{amsart} preamble that restates the article's own declarations and macros used by the retained lines, namely the environments \texttt{proposition} on separate counters, together with the standard packages those lines need.  The retained environments print in this document as Proposition 1 in order of appearance, whereas the article prints the same items with the numbering of its own full document.  The complete native source of the article as distributed in the arXiv e-print for this exact version is bundled unchanged as \texttt{sources/197017/source0.tex}.
\begin{proposition}\label{prop:lowerbound2}
For $p\ge 2$, $C(p,2)\ge p-1$.
\end{proposition}

\begin{proof}
Define the vector field $u_k\colon \R^2 \to \R^2$ as 
\begin{equation*}
u_k(x) = Q x \ln(\abs{x})^k \mathds{1}_{\abs{x}<1},
\end{equation*}
where $x=(x_1, x_2) \in \R^2$,
$\abs{x} = \sqrt{x_1^2 + x_2^2}$,
$Q = \mat{0 & 1 \\ -1 & 0}$, $k\in \mathbb N,\ k\ge 1$,
and $\mathds{1}_{\abs{x}<1}$
is the indicator function of the open unit disk
$B_{\R^2}(0,1)$.
Let $r = \abs{x}$. Then
\begin{equation*}
\nabla u_k = \ln(r)^k Q + k \ln(r)^{k-1} r^{-2} M(x),
\end{equation*} where 
\begin{equation*}
M(x) = \mat{x_1x_2 & x_2^2 \\ -x_1^2 & -x_1x_2}.
\end{equation*}
It follows that
\begin{align*}
\mathcal E(u_k)(x)
= k \ln(r)^{k-1} r^{-2} C(x) \mathds{1}_{\abs{x}<1},
\end{align*}
where 
\begin{equation*}
C(x) = \frac{1}{2}(M(x)+M(x)^T)
= \mat{x_1x_2 & (x_2^2-x_1^2)/2
\\ (x_2^2-x_1^2)/2 & -x_1x_2}
\end{equation*}
and that
\begin{equation*}
\mathcal A(u_k)(x) =
\left(\ln(r)^k + \frac{k}{2} \ln(r)^{k-1}\right)
Q \mathds{1}_{\abs{x}<1}
\end{equation*}
Taking the Frobenius norm, since
$|C(x)| = r^2 / \sqrt{2}$ 
it follows that (writing $t=-\ln(r)$)
\begin{align*}
|\mathcal E(u_k)(x)|
&= \abs{k \ln(r)^{k-1} r^{-2}} |C(x)|
= \abs{k} \abs{\ln r}^{k-1} r^{-2} (r^2 / \sqrt{2})
= \frac{\abs{k}}{\sqrt{2}} t^{k-1}\\
|\mathcal A(u_k)(x)|
&= \abs{\ln(r)^{k-1}(\ln r + k/2)} |Q|
= \sqrt{2} \abs{\ln r}^{k-1} \abs{\ln r + k/2}\\
&= \sqrt{2} t^{k-1} \abs{-t + k/2}
= \sqrt{2} t^{k-1} \abs{t - k/2}.
\end{align*}
Taking $L^p$ norms we obtain
\begin{align*}
\norm{\mathcal E(u_k)}_{L^p}^p& = \int_{B_{\R^2}(0,1)}
|\mathcal E(u_k)(x)|^p \diff x
= \int_0^{2\pi} \int_0^1
\left(\frac{\abs{k}}{\sqrt{2}} t^{k-1}\right)^p
r \diff r \diff \theta\\&
= 2\pi \left(\frac{\abs{k}}{\sqrt{2}}\right)^p
\int_0^1 r (-\ln r)^{p(k-1)} \diff r.
\end{align*}
The integral can also be written as
\begin{equation*}
\int_0^\infty e^{-2t} t^{p(k-1)} dt =
\left(\frac{1}{2}\right)^{p(k-1)+1}
\Gamma(p(k-1)+1)
=\mathrel{:}\left(\frac{1}{2}\right)^{p(k-1)+1}I_p.
\end{equation*}
Thus it follows that
\begin{align*}
\norm{\mathcal E(u_k)}_{L^p}^p
&= 2\pi k^p 2^{-p/2} 2^{-(p(k-1)+1)}
\Gamma(p(k-1)+1)\\
&=2\pi k^p 2^{-p/2} 2^{-(p(k-1)+1)} I_p.
\end{align*}
To calculate the norm of the anti-symmetric part,
first observe that 
\begin{align*}
\norm{\mathcal A(u_k)}_{L^p}^p& =
\int_{B_{\R^2}(0,1)} \abs{\mathcal A(u_k)(x)}^p dx
= \int_0^{2\pi} \int_0^1
\left(
\sqrt{2} t^{k-1} \abs{t - k/2}
\right)^p r dr d\theta\\
&= 2\pi (\sqrt{2})^p\int_0^\infty e^{-2t} t^{p(k-1)}
\abs{t - k/2}^p \diff t.
\end{align*}
Define 
\begin{equation*}
J_p\mathrel{:}=\int_0^\infty e^{-t} t^{p(k-1)}
\left|\frac tk - 1\right|^p\diff t
\end{equation*}
so that as a simple calculation via the
change of variable $s=2t$ shows we have
\begin{equation*}
\norm{\mathcal A(u_k)}_{L^p}^p
=2\pi 2^{p/2} J_p 2^{-p(k-1)-1-p}k^p
\end{equation*}
We now want to find a lower bound for the function
$f_k(p) \mathrel{:}=
\norm{\mathcal A(u_k)}_{L^p}/
\norm{\mathcal E(u_k)}_{L^p}$.
Thanks to our calculations we have
\begin{align}
f_k(p)^p = \frac{\norm{\mathcal A(u_k)}_{L^p}^p}
{\norm{\mathcal E(u_k)}_{L^p}^p}
&= \frac{2\pi (\sqrt{2})^p k^p2^{-p} J_p}
{2\pi (k/\sqrt{2})^p I_p}
\nonumber \\
&= \frac{J_p}{I_p} \label{eq:fk_ratio} 
\end{align}
We now prove that $\lim_{k\to \infty} J_p/I_p\ge (p-1)^p$,
which will conclude the proof.
First notice that the ratio admits
an interpretation as an expectation: 
define a probability measure
$d\mu(t)$ on $(0, \infty)$ by
\begin{equation}
 d\mu(t) = \frac{e^{-t} t^{p(k-1)}}{I_p} dt. 
\end{equation}
Then,
\begin{equation}
\frac{J_p}{I_p}
= \int_0^\infty \left|\frac tk - 1\right|^p
\frac{e^{-t} t^{p(k-1)}}{I_p} dt
= \mathbb{E}[\abs{X_k - 1}^p] 
\end{equation}
where $X_k$ is a random variable that is
$(p(k-1)+1,1/k)$ Gamma-distributed.
By Jensen's inequality, $\mathbb{E}[\phi(X)]
\ge \phi(\mathbb{E}[X])$
for convex functions $\phi$, hence
\begin{equation}
\mathbb{E}[|X_k-1|^p] \ge \abs{\mathbb{E}[X_k] - 1}^p 
\end{equation}
Some straightforward calculations
with the $\Gamma$ function prove that
\begin{align*}
\mathbb{E}[X_k]= \frac{(p(k-1)+1)}k,
\end{align*}
so
\begin{align*}
\frac{J_p}{I_p}& \ge \abs{\mathbb{E}[X_k] - 1}^p \\
&= \left|p\frac{k-1}k+\frac1k-1\right|^p
\end{align*}
Substituting this lower bound for $J_p/I_p$ back
into~\eqref{eq:fk_ratio} gives us
\begin{equation*}
 f_k(p) \ge (p-1)\frac{k-1}{k}.
\end{equation*}
Since $C(p,2)\ge f_k(p)$ for any $k$ it follows that
\begin{equation*}
    C(p,2)\ge \liminf_{k\to \infty}f_k(p)=p-1.
\end{equation*}
\end{proof}

\end{document}
